% sections/data.tex: Section 3. Counts: results/family_counts.csv, results/summary_stats_families.csv (E19), results/panel_windows_auto.csv.
\section{Data}\label{sec:data}

Three samples produce the estimates. The first is a one-percent consumer panel of U.S.\ credit files, observed semiannually from 2005 to 2009 and monthly from January 2010 to December 2025, which holds, across its six loan families (auto, first mortgage, home equity, personal installment, student and credit card), 67.2 million loans, more than 3.8 million borrowers and 5.8 million first ninety-day delinquencies; its 9.7 million auto loans, of which 7.6 million form the analysis sample, are the main panel. The second is a ten-percent trade archive, observed at 57 archive months from 2005 to April 2024 with a monthly score file from 2010, on which the within-borrower designs run: of its 15.5 million borrowers with two or more auto loans (83.5 million loans), the 13.5 million whose loans meet the contract restrictions below, 65 million loans. The third is the covariate-matched ten-percent extract, the loans of the one-percent consumers nested in the archive, which alone has age and income; the specification search of Appendix \ref{appendix:bridge}, the deficiency contrast and the term margin by score decile run on it (72,906 lifetime defaults). The population estimates run on the one-percent panel rather than the archive because the panel is monthly and records income and age on every loan, which the archive's quarterly snapshots do not; the archive's size is what the within-borrower designs need, since only a fraction of borrowers hold two auto loans with a default in the window. Appendix \ref{appendix:data} records the files, frequencies and restrictions of every exhibit.
\begin{table}[H]
\centering
\caption{Summary statistics by loan family}
\label{tab:summary_families}
\small
\resizebox{\textwidth}{!}{\begin{tabular}{lrrrrrr}
\toprule
 & Auto & Credit card & First mortgage & Student & Personal installment & Home equity \\
\midrule
\multicolumn{7}{l}{\emph{Panel A. Coverage: the one-percent consumer panel, 2005--2025}} \\
Loans (millions) & 9.66 & 34.20 & 6.12 & 8.61 & 8.08 & 0.53 \\
Borrowers (millions) & 2.20 & 3.81 & 1.46 & 0.94 & 1.53 & 0.28 \\
Lenders & 18,690 &  9,758 & 15,188 &  5,699 & 30,273 &  7,527 \\
Loans per borrower & 4.4 & 9.0 & 4.2 & 9.2 & 5.3 & 1.9 \\
Borrowers with two or more loans (percent) & 73.9 & 75.9 & 78.5 & 85.8 & 65.9 & 46.2 \\
Opening year, 10th to 90th percentile & 2003--2022 & 1995--2022 & 2001--2021 & 2001--2020 & 2002--2023 & 2002--2021 \\
Loans with a score and an income (percent) & 98.4 & 95.2 & 99.3 & 97.9 & 98.3 & 99.3 \\
\addlinespace
\multicolumn{7}{l}{\emph{Panel B. The contract at origination (medians)}} \\
Balance or limit (\$) &  19,089 &   1,215 & 169,000 &   3,744 &   2,791 &  36,366 \\
Scheduled monthly payment (\$) &   387 &    29 & 1,263 &    50 &   132 &   359 \\
Contract length (months) &  60 & revolving & 360 & 120 &  24 & 180 \\
Payment to monthly income & 0.114 & 0.010 & 0.290 & 0.021 & 0.053 & 0.088 \\
\addlinespace
\multicolumn{7}{l}{\emph{Panel C. The borrower at first sight}} \\
Credit score, median & 698 & 724 & 757 & 653 & 621 & 734 \\
Credit score, 10th to 90th percentile & 550--807 & 565--817 & 605--822 & 488--769 & 492--770 & 576--817 \\
Subprime, score below 620 (percent) & 25.0 & 19.7 & 12.0 & 38.2 & 49.6 & 16.4 \\
Modeled monthly income, median (\$) & 3,417 & 3,417 & 4,167 & 2,500 & 2,583 & 4,083 \\
Income below \$2,500 a month (percent) & 24.0 & 22.1 & 7.2 & 49.1 & 43.7 & 6.8 \\
Age, median & 44 & 48 & 47 & 31 & 46 & 47 \\
\addlinespace
\multicolumn{7}{l}{\emph{Panel D. Outcomes over the life of the loan}} \\
Ever 30 days past due (percent) & 12.5 & 13.9 & 12.1 & 17.1 & 11.3 & 13.9 \\
Ever 90 days past due (percent) & 3.5 & 9.3 & 6.4 & 15.9 & 5.8 & 8.9 \\
Charged off (percent) & 4.9 & 9.7 & 0.2 & 0.8 & 6.3 & 3.1 \\
90 days past due within 12 months (percent) & 0.6 & 1.9 & 0.2 & 0.5 & 1.9 & 0.3 \\
90 days past due within 24 months (percent) & 1.5 & 3.6 & 0.9 & 2.2 & 3.5 & 1.6 \\
Months to first 90 days past due, median & 28 & 32 & 55 & 61 & 13 & 49 \\
Charge-off amount over balance, median & 0.49 & 1.88 & 0.85 & 1.23 & 0.86 & 0.97 \\
\addlinespace
\multicolumn{7}{l}{\emph{Panel E. The ten-percent trade archive, 2005--2024 (within-borrower designs)}} \\
Loans (millions) & 89.0 &  & 58.2 &  & 73.4 & 4.9 \\
Borrowers (millions) & 21.1 &  & 14.1 &  & 14.4 & 2.6 \\
Borrowers with two or more loans (millions) & 15.5 &  & 11.1 &  & 9.4 & 1.2 \\
Their loans (millions) & 83.5 &  & 55.1 &  & 68.4 & 3.5 \\
Charged off (percent) & 4.1 &  & 0.1 &  & 5.7 & 2.5 \\
\bottomrule
\end{tabular}
}
\vspace{0.5em}
\begin{minipage}{0.96\textwidth}
\noindent \small \textit{Notes:} Panels A to D describe the one-percent consumer panel, semiannual from 2005 to 2009 and monthly from 2010 to 2025, with every loan of each family that the archive reports. A borrower is a hashed consumer identity; a lender is a hashed furnisher identity. The balance is the original balance for installment loans and the high credit at first sight for cards; the scheduled payment is the payment at first sight, for cards the minimum; a card's contract has no length. Modeled income is the bureau's estimate of monthly income at the quarter of first sight; the score is the bureau's risk score at the month of first sight; the row ``loans with a score and an income'' is the share of loans on which both covariates are recorded, the coverage of every exhibit that conditions on them. Delinquency and charge-off are the first occurrence over the loan's observed life; the twelve- and twenty-four-month rates condition on loans opened at least that long before the panel's last month and, for installment loans, on contracts of at least that length. The charge-off amount is divided by the balance at origination, so that it exceeds one where interest and fees were capitalized. Panel E describes the ten-percent trade archive, 57 snapshots from 2005 to 2024, on which the within-borrower designs run; cards and student loans have no extract from it.
\end{minipage}
% replication: Code/extract/E19_summary_stats.R, Code/exhibits/make_table1.R
\end{table}

The credit-bureau panel records every loan type a household holds, and the loan types differ in the three respects that matter for measuring a discount rate: the contract, which fixes the payment path and the horizon; the population, which fixes whose preferences the estimate describes; and the instruments, which fix what can be identified. Table \ref{tab:loan_types} in Appendix \ref{appendix:data} sets the loan types side by side. Auto loans and credit cards are the laboratories for the low end of the income distribution, and the auto-loan panel is the main sample for the estimates that follow. First mortgages are the laboratory for long horizons, up to thirty years, and the one whose instruments, house prices by origination cohort and state recourse law, are the strongest available; Section \ref{sec:equity} reports what they establish. Student loans have the one payment fixed by statute, the resumption of October 2023; Appendix \ref{appendix:timing} reports that design and the adjustable-rate resets of first mortgages in full. The panel is a consumer credit panel drawn from the files of a national credit bureau, one of the panels surveyed by \citet{gibbs_consumer_2025}, and every measure follows their guidance except where Appendix \ref{appendix:data} states otherwise.

\paragraph{The auto-loan panel.} Auto loans offer an attractive empirical setting: maturities vary substantially, payments are highly salient, and auto loans are one of the most common forms of installment credit, with wide participation across the income distribution. The analysis sample restricts the panel's 9.7 million auto loans to contracts of 12 to 96 months with a positive scheduled payment, a score and an income, and an identified lender: 7.6 million loans at 15,532 lenders. The typical repayment horizon is five years, and the three most popular maturities, 36, 60 and 72 months, account for over half of the loans; origination volumes are steady from 2010 onward (Figure \ref{fig:origination_volume}, Appendix \ref{appendix:figures}). Table \ref{tab:summary_families} summarizes the key variables by family, Table \ref{tab:loan_summary} in Appendix \ref{appendix:procedures} the auto panel's quartiles, and Appendix \ref{appendix:procedures} the construction.
